\documentclass[11pt,a4paper]{article}

\usepackage[utf8]{inputenc}
\usepackage[T1]{fontenc}
\usepackage[margin=2.5cm]{geometry}
\usepackage{hyperref}
\usepackage{xcolor}
\usepackage{listings}
\usepackage{enumitem}
\usepackage{booktabs}
\usepackage{longtable}
\usepackage{amsmath,amssymb}
\usepackage{fancyhdr}
\usepackage{titlesec}

\definecolor{isls-blue}{HTML}{1a5276}
\definecolor{isls-orange}{HTML}{e67e22}
\definecolor{isls-green}{HTML}{27ae60}
\definecolor{isls-red}{HTML}{c0392b}
\definecolor{isls-gray}{HTML}{7f8c8d}
\definecolor{code-bg}{HTML}{f7f9fb}

\lstdefinestyle{rust}{
  language=C,
  basicstyle=\ttfamily\small,
  keywordstyle=\color{isls-blue}\bfseries,
  commentstyle=\color{isls-gray}\itshape,
  stringstyle=\color{isls-green},
  backgroundcolor=\color{code-bg},
  frame=single,
  rulecolor=\color{isls-gray!30},
  breaklines=true,
  showstringspaces=false,
  tabsize=4,
  morekeywords={fn,let,mut,pub,struct,impl,enum,use,mod,self,Self,
    crate,super,where,trait,for,in,if,else,match,return,async,await,
    Ok,Err,Some,None,Vec,String,HashMap,BTreeMap,BTreeSet,
    Result,Option,PathBuf,VecDeque,
    f64,f32,usize,bool,u32,u64,i64,Box,dyn,Arc,Mutex},
}
\lstdefinestyle{bash}{
  language=bash,
  basicstyle=\ttfamily\small,
  backgroundcolor=\color{code-bg},
  frame=single,
  rulecolor=\color{isls-gray!30},
  breaklines=true,
  showstringspaces=false,
}

\pagestyle{fancy}
\fancyhf{}
\fancyhead[L]{\small\textcolor{isls-gray}{ISLS I5 Spec --- Sensorium}}
\fancyhead[R]{\small\textcolor{isls-gray}{v1.0.0 --- \today}}
\fancyfoot[C]{\thepage}

\titleformat{\section}
  {\Large\bfseries\color{isls-blue}}{\thesection}{1em}{}
\titleformat{\subsection}
  {\large\bfseries\color{isls-blue!80}}{\thesubsection}{1em}{}
\titleformat{\subsubsection}
  {\normalsize\bfseries\color{isls-blue!60}}{\thesubsubsection}{1em}{}

\hypersetup{
  colorlinks=true,
  linkcolor=isls-blue,
  urlcolor=isls-orange,
}

\begin{document}

\begin{titlepage}
\centering
\vspace*{2.5cm}

{\Huge\bfseries\color{isls-blue} ISLS I5 Specification}\\[0.5cm]
{\LARGE\color{isls-orange} Sensorium}\\[0.3cm]
{\large\color{isls-gray} Infogenetik-Dashboard, SGB-Tracking,\\
Zeitreihen-Visualisierung, Autonome Zyklen}\\[2cm]

{\large
\begin{tabular}{rl}
\textbf{System:}       & ISLS --- Intelligent Semantic Ledger Substrate \\
\textbf{Version:}      & v6.4 (I5) \\
\textbf{Author:}       & Sebastian Klemm \\
\textbf{Date:}         & \today \\
\textbf{Status:}       & Specification \\
\textbf{Depends:}      & I4 (complete), GUI-Wiring (complete) \\
\textbf{Branch:}       & \texttt{claude/implement-isls-i5-spec} \\
\textbf{Regel 13:}     & Jedes Feature MUSS im Studio erreichbar sein. \\
                        & Was nur in der CLI existiert ist nicht fertig. \\
\end{tabular}
}

\vfill

{\small\color{isls-gray}
Mithras Substructure University\\
Repository: \texttt{https://github.com/LashSesh/graphity}\\
License: MIT
}
\end{titlepage}

\tableofcontents
\newpage


% ══════════════════════════════════════════════════════════════════════════════
\section{Executive Summary}
% ══════════════════════════════════════════════════════════════════════════════

ISLS hat jetzt: Generierung (D1--D8), Konsens (M1), Studio (S1), 
Messung (I1), Evolution (I2), Metamorphose (I3), Cross-Language (I4), 
und eine laufende \"Olbohrinsel die 12+ Auto-Norms aus 150+ Repos 
kristallisiert hat. Alles funktioniert --- aber der Operator 
sieht nur Zahlen. ``Auto-Norms: 12''. Was \emph{bedeutet} das? 
Wie hat sich die Norm-Dichte \"uber die letzten 24 Stunden 
entwickelt? Welche Gene sind stabil? Wo liegt die SGB? Konvergiert 
die Fitness? Welche Dom\"anen sind ges\"attigt, welche hungern?

I5 installiert das \textbf{Sensorium}: ein visuelles Dashboard 
im Studio das den Zustand des Infogenetik-Systems in Echtzeit 
zeigt --- nicht als Zahlenkolonnen sondern als Graphen, 
Heatmaps und Trajektorien. Plus: autonome Metamorphose-Zyklen 
die das System ohne menschlichen Eingriff durch periodische 
Evolve-Zyklen verbessern.

\subsection{Was I5 Produziert}

\begin{enumerate}[nosep]
\item \textbf{Vierter Studio-Modus: Sensorium ($\blacklozenge$):} 
      Ein Dashboard mit Zeitreihen und Visualisierungen. 
      Kein Chat. Nur Daten.
\item \textbf{Zeitreihen-Persistenz:} Norm-Stats, Fitness, SGB, 
      Scrape-Raten werden periodisch in eine 
      \texttt{timeseries.jsonl} geschrieben.
\item \textbf{Sechs Visualisierungen:} SGB-Trajektorie, 
      Fitness-Verteilung, Norm-Wachstum, Scrape-Rate, 
      Codematrix-Radar, Gen-Topologie.
\item \textbf{Autonomer Evolve-Zyklus:} Optionaler Cronjob 
      der w\"ochentlich \texttt{isls evolve --from . --delta 
      "optimize"} ausf\"uhrt und das System sich selbst 
      neu drucken l\"asst.
\end{enumerate}


% ══════════════════════════════════════════════════════════════════════════════
\section{Architektur-Regel 13}
% ══════════════════════════════════════════════════════════════════════════════

\begin{center}
\fbox{\parbox{0.85\textwidth}{
\textbf{Regel 13 (Studio-Pflicht):} Jedes Feature MUSS im 
Studio erreichbar sein. Ein Feature das nur als CLI-Befehl 
oder API-Endpoint existiert ist NICHT FERTIG. Die CLI ist 
Fallback, nicht Prim\"arinterface. Diese Regel gilt r\"uckwirkend 
und f\"ur alle zuk\"unftigen Specs.
}}
\end{center}

Konsequenz f\"ur I5: Alles was hier spezifiziert wird, hat 
eine Studio-Komponente. Kein W-Paket ohne GUI.


% ══════════════════════════════════════════════════════════════════════════════
\section{W1: Zeitreihen-Persistenz}
\label{sec:w1}
% ══════════════════════════════════════════════════════════════════════════════

\subsection{Problem}

Die aktuellen Daten (\texttt{norms.json}, \texttt{fitness.json}, 
\texttt{metrics.jsonl}) speichern den \emph{aktuellen Zustand} 
aber nicht die \emph{Geschichte}. Man sieht ``SGB: 0.87'' aber 
nicht ``SGB war vor 24h bei 0.82 und steigt''.

\subsection{Timeseries Snapshot}

Alle 15 Minuten schreibt ISLS einen Snapshot in 
\texttt{\textasciitilde/.isls/timeseries.jsonl}:

\begin{lstlisting}[style=rust]
#[derive(Debug, Serialize, Deserialize)]
pub struct TimeseriesEntry {
    pub timestamp: String,       // ISO 8601
    
    // Norm-System
    pub norms_builtin: usize,
    pub norms_auto: usize,
    pub norms_injected: usize,
    pub norms_total: usize,
    pub candidates_total: usize,
    pub candidates_observing: usize,
    pub candidates_promoted: usize,
    pub observations_total: usize,
    pub domains_total: usize,
    
    // Fitness
    pub fitness_mean: f64,       // Durchschnitt aller Norm-Fitnesswerte
    pub fitness_min: f64,
    pub fitness_max: f64,
    pub fitness_count: usize,    // Anzahl getrackter Normen
    
    // SGB (aus letzter Generierung, falls vorhanden)
    pub sgb: Option<f64>,
    pub last_compile_success: Option<bool>,
    pub last_codematrix_resonance: Option<f64>,
    
    // Scraping
    pub repos_scraped_total: usize,
    pub keywords_active: usize,
    pub keywords_suggested: usize,
    
    // Genom
    pub genes_count: usize,
    pub cross_language_norms: usize,
}
\end{lstlisting}

\subsection{Snapshot-Trigger}

In \texttt{isls-gateway}: ein \texttt{tokio::spawn} Background-Task 
der alle 15 Minuten einen Snapshot schreibt:

\begin{lstlisting}[style=rust]
async fn timeseries_recorder(state: AppState) {
    let mut interval = tokio::time::interval(
        Duration::from_secs(900) // 15 min
    );
    loop {
        interval.tick().await;
        if let Ok(entry) = collect_timeseries_snapshot(&state) {
            append_timeseries_entry(&entry);
        }
    }
}
\end{lstlisting}

\subsection{Datei-Format}

Ein JSON-Objekt pro Zeile (JSONL). Append-only. Keine Rotation 
(eine Zeile pro 15 Minuten $=$ 96/Tag $=$ $\sim$35K/Jahr $=$ 
$\sim$10~MB/Jahr). Kein Aufr\"aumen n\"otig.

\subsection{API Endpoint}

\begin{lstlisting}[style=bash]
GET /api/timeseries?hours=24
GET /api/timeseries?hours=168  (1 Woche)
GET /api/timeseries?hours=720  (1 Monat)

Response: { "entries": [...], "count": 96 }
\end{lstlisting}

Liest \texttt{timeseries.jsonl}, filtert nach Zeitfenster, 
gibt die Eintr\"age zur\"uck.


% ══════════════════════════════════════════════════════════════════════════════
\section{W2: Sensorium-Modus im Studio}
\label{sec:w2}
% ══════════════════════════════════════════════════════════════════════════════

\subsection{Vierter Modus}

Ein neues Icon im Mode-Strip: $\blacklozenge$ (oder $\diamondsuit$ 
oder ein Auge-Symbol). Klick wechselt zum Sensorium. Kein Chat. 
Kein Input-Feld. Nur Visualisierungen.

\subsection{Layout}

Sechs Karten in einem responsiven Grid (2 Spalten Desktop, 
1 Spalte Mobile). Jede Karte hat einen Titel und einen 
Graphen oder eine Visualisierung. Alle Daten kommen von 
\texttt{GET /api/timeseries} und den bestehenden Norm/Fitness/
Metrics-Endpoints.

\begin{lstlisting}[style=bash]
+--+--------------------------------------------------+
|  |                                                    |
|  |  +---------------------+  +---------------------+ |
| *|  | SGB-Trajektorie     |  | Norm-Wachstum       | |
| o|  | [Liniengraph]       |  | [Liniengraph]       | |
| .|  +---------------------+  +---------------------+ |
| .|  +---------------------+  +---------------------+ |
|  |  | Fitness-Verteilung  |  | Scrape-Rate         | |
|  |  | [Balkengraph]       |  | [Liniengraph]       | |
|  |  +---------------------+  +---------------------+ |
|  |  +---------------------+  +---------------------+ |
|  |  | Codematrix-Radar    |  | Gen-Topologie       | |
|  |  | [Radardiagramm]     |  | [Cluster-Map]       | |
|  |  +---------------------+  +---------------------+ |
|  |                                                    |
+--+--------------------------------------------------+
\end{lstlisting}

Kein Input-Feld im Sensorium. Keine Nachrichten. Nur 
Visualisierungen die sich alle 30 Sekunden aktualisieren.

\subsection{Karte 1: SGB-Trajektorie}

\textbf{Datenquelle:} \texttt{timeseries.jsonl} $\to$ 
\texttt{sgb} Feld \"uber Zeit.

\textbf{Darstellung:} Liniengraph. X-Achse: Zeit (letzte 
24h/7d/30d, umschaltbar). Y-Achse: SGB (0.0--1.0). 
Die Linie zeigt wie der strukturelle Anteil w\"achst. 
Eine horizontale Linie bei 1.0 markiert das Ziel 
(``Full Understanding'').

\textbf{Farbe:} \texttt{--accent} (orange) f\"ur die Linie. 
\texttt{--text-dim} f\"ur die Ziellinie.

\textbf{Rendering:} Inline SVG. Kein Chart-Framework. 
Einfache \texttt{<polyline>} im \texttt{<svg>}. Die 
Datenpunkte werden aus der Timeseries berechnet.

\subsection{Karte 2: Norm-Wachstum}

\textbf{Datenquelle:} \texttt{timeseries.jsonl} $\to$ 
\texttt{norms\_total}, \texttt{norms\_auto}, 
\texttt{candidates\_total}.

\textbf{Darstellung:} Drei Linien \"ubereinander:
\begin{itemize}[nosep]
\item Builtin (flach, \"andert sich nie) in \texttt{--text-dim}
\item Auto-Norms (steigend) in \texttt{--accent-success} (gr\"un)
\item Candidates (schnell steigend) in \texttt{--accent-info} (blau)
\end{itemize}

\subsection{Karte 3: Fitness-Verteilung}

\textbf{Datenquelle:} \texttt{GET /api/norms/fitness}

\textbf{Darstellung:} Horizontales Balkendiagramm. Jede Norm 
ein Balken. Sortiert nach Fitness (h\"ochste oben). 
Farbe: gr\"un ($\phi > 0.8$), orange ($0.5$--$0.8$), 
rot ($< 0.5$). Max 20 Normen anzeigen. Wenn mehr: 
``+X weitere'' Text.

\subsection{Karte 4: Scrape-Rate}

\textbf{Datenquelle:} \texttt{timeseries.jsonl} $\to$ 
\texttt{repos\_scraped\_total} (Differenz zwischen 
aufeinanderfolgenden Eintr\"agen $=$ Rate).

\textbf{Darstellung:} Liniengraph. Y-Achse: Repos/Stunde. 
Zeigt ob die \"Olbohrinsel gleichm\"a{\ss}ig pumpt oder 
stockt.

\subsection{Karte 5: Codematrix-Radar}

\textbf{Datenquelle:} \texttt{GET /api/timeseries} $\to$ 
letzter Eintrag mit \texttt{last\_codematrix\_resonance}, 
oder \texttt{GET /api/metrics} $\to$ letzter Forge-Run 
mit Codematrix-Daten.

\textbf{Darstellung:} Radardiagramm (Pentagon) mit den 
5 Achsen R, F, T, S, E. Zeigt das durchschnittliche 
Codematrix-Profil der letzten Generierung. Wenn keine 
Generierung vorliegt: ``Kein Forge-Run vorhanden.''

\textbf{Rendering:} SVG Polygon. 5 Achsen als Linien vom 
Zentrum. Datenpunkte als gef\"ulltes Polygon in 
\texttt{--accent} mit 30\% Opacity.

\subsection{Karte 6: Gen-Topologie}

\textbf{Datenquelle:} \texttt{GET /api/norms/genome}

\textbf{Darstellung:} Cluster-Visualisierung. Jedes Gen 
ist ein Kreis. Gr\"o{\ss}e $\propto$ Anzahl der Normen im 
Gen. Farbe: \texttt{--accent-success} f\"ur hohe Fitness, 
\texttt{--accent-error} f\"ur niedrige. Gen-Name als Label. 
Singletons kleiner und in \texttt{--text-dim}.

Wenn zu wenig Daten (< 10 Metriken-Eintr\"age): 
``Nicht genug Daten f\"ur Gen-Clustering.''

\subsection{Zeitfenster-Umschaltung}

Oben rechts im Sensorium: drei Buttons (keine Tabs, 
nur Text-Links):

\textbf{24h} | \textbf{7d} | \textbf{30d}

Klick l\"adt die Timeseries mit dem entsprechenden 
\texttt{hours}-Parameter und aktualisiert alle Graphen.
Default: 24h.


% ══════════════════════════════════════════════════════════════════════════════
\section{W3: SVG-Rendering-Engine}
\label{sec:w3}
% ══════════════════════════════════════════════════════════════════════════════

Kein externes Chart-Framework. Alle Graphen werden als 
inline SVG im JavaScript erzeugt. Eine kleine 
Rendering-Bibliothek innerhalb der \texttt{studio.html}:

\begin{lstlisting}[style=bash]
function renderLineChart(container, data, options) {
    // data: [{x: timestamp, y: value}, ...]
    // options: {width, height, color, yMin, yMax, 
    //           xLabel, yLabel, targetLine}
    // Erzeugt <svg> mit:
    //   - Achsen (--border Farbe)
    //   - Gitternetzlinien (--text-dim, gestrichelt)
    //   - Datenlinie (<polyline>, options.color)
    //   - Achsenbeschriftungen (--text-secondary)
    //   - Optional: targetLine (horizontale Linie)
}

function renderBarChart(container, data, options) {
    // data: [{label, value, color}, ...]
    // Horizontale Balken, sortiert
}

function renderRadarChart(container, axes, values, options) {
    // axes: ["R", "F", "T", "S", "E"]
    // values: [0.62, 0.78, 0.91, 0.88, 0.55]
    // Pentagon mit gefuelltem Polygon
}

function renderClusterMap(container, clusters, options) {
    // clusters: [{name, size, fitness, norms}, ...]
    // Kreise mit Groesse und Farbe
}
\end{lstlisting}

Alle Funktionen erzeugen reines SVG. Kein Canvas. Kein 
WebGL. Kein D3.js. Einfaches DOM-Building mit 
\texttt{document.createElementNS('http://www.w3.org/2000/svg', 
...)}.

\subsection{Responsive}

SVG nutzt \texttt{viewBox} f\"ur automatische Skalierung. 
Karten-Container nutzen CSS Grid mit \texttt{minmax(300px, 1fr)}. 
Auf Mobile: eine Spalte, Karten stacken vertikal.


% ══════════════════════════════════════════════════════════════════════════════
\section{W4: Autonomer Evolve-Zyklus}
\label{sec:w4}
% ══════════════════════════════════════════════════════════════════════════════

\subsection{Problem}

Das System verbessert sich durch Norm-Akkumulation (Scraping) 
und durch Generierungserfahrung (Fitness). Aber die 
Generierungen passieren nur wenn der User aktiv etwas forgt. 
Ohne User-Aktivit\"at lernt das System nur durch Scraping, 
nicht durch eigene Emissionen.

\subsection{L\"osung: Periodischer Self-Forge}

Ein optionaler Hintergrund-Task der periodisch eine 
Standard-App generiert und die Emission beobachtet. 
Nicht um die App zu benutzen --- nur um Fitness-Daten 
zu sammeln und die Norm-Qualit\"at zu testen.

\begin{lstlisting}[style=rust]
async fn auto_evolve_cycle(state: AppState) {
    let descriptions = vec![
        "Pet shop with animals, owners, and appointments",
        "Hotel booking system with rooms, guests, and reservations",
        "Blog platform with posts, authors, comments, and tags",
        "Task manager with projects, tasks, and team members",
        "Inventory system with products, warehouses, and orders",
    ];
    
    let mut interval = tokio::time::interval(
        Duration::from_secs(6 * 3600)  // alle 6 Stunden
    );
    
    loop {
        interval.tick().await;
        
        // Zufaellige Beschreibung waehlen
        let desc = descriptions.choose(&mut rng).unwrap();
        
        tracing::info!("[AutoEvolve] Starting: {}", desc);
        
        let output = temp_dir().join(format!(
            "auto-evolve-{}", timestamp()
        ));
        
        match forge_chat(desc, &output, oracle.clone()).await {
            Ok(result) => {
                tracing::info!(
                    "[AutoEvolve] Success: {} files, {} LOC, {:.0}s",
                    result.files, result.loc, result.duration
                );
                // Scrape das Ergebnis fuer Selbstbeobachtung
                scrape_path(&output).ok();
                // Aufraeumen
                std::fs::remove_dir_all(&output).ok();
            }
            Err(e) => {
                tracing::warn!("[AutoEvolve] Failed: {}", e);
            }
        }
    }
}
\end{lstlisting}

\subsection{Aktivierung}

Der Auto-Evolve-Zyklus ist \textbf{standardm\"a{\ss}ig deaktiviert}. 
Aktivierung \"uber:

\begin{itemize}[nosep]
\item CLI-Flag: \texttt{isls serve --port 8420 --ollama --auto-evolve}
\item Service-Datei: 
      \texttt{ExecStart=... --auto-evolve}
\item Studio: Toggle im Sensorium-Modus
\end{itemize}

\subsection{Studio-Toggle}

Im Sensorium-Modus: ein Toggle-Schalter oben rechts neben 
den Zeitfenster-Buttons:

\begin{lstlisting}[style=bash]
  24h | 7d | 30d          [Auto-Evolve: AUS]
\end{lstlisting}

Klick wechselt zwischen AUS und AN. Status wird in 
\texttt{AppState} gespeichert (nicht persistent \"uber 
Reboot --- muss bewusst aktiviert werden).

Wenn AN: der Sensorium zeigt zus\"atzlich die 
Auto-Evolve-History:

\begin{lstlisting}[style=bash]
  Letzter Auto-Evolve: vor 3h
    "Hotel booking system" - 45 files, 162s, PASS
  Naechster Auto-Evolve: in 3h
  Heute: 3/4 PASS (75%)
\end{lstlisting}

\subsection{Selbstbeobachtung}

Jede Auto-Evolve-Emission wird gescrapt (\texttt{scrape\_path}). 
Dadurch:
\begin{itemize}[nosep]
\item Fitness wird aktualisiert (I1/I2)
\item Neue Observations entstehen
\item Codematrix-Daten werden gesammelt
\item Timeseries wird gef\"uttert
\end{itemize}

Das System generiert, beobachtet sich selbst, und lernt --- 
ohne menschlichen Eingriff. Der Feedback-Loop dreht autonom.


% ══════════════════════════════════════════════════════════════════════════════
\section{W5: SGB-Tracking und Konvergenz-Alarm}
\label{sec:w5}
% ══════════════════════════════════════════════════════════════════════════════

\subsection{SGB-Berechnung}

Die SGB (Structural-Generative Boundary) wird bei jeder 
Generierung berechnet:

\begin{equation}
\text{SGB} = \frac{k_S}{k_S + k_O}
\end{equation}

Aktuell: $\text{SGB} \approx 0.87$ (44 strukturelle / 50 
total). Dieser Wert wird in \texttt{GenerationMetrics} und 
in der Timeseries gespeichert.

\subsection{Konvergenz-Alarm}

Wenn die SGB \"uber 3 aufeinanderfolgende Generierungen 
\emph{sinkt}: Warnung im Sensorium.

\begin{lstlisting}[style=bash]
  ! SGB sinkt: 0.89 -> 0.87 -> 0.85 (3 Runs)
    Moegliche Ursache: Neue Normen mit niedrigem 
    Fitness-Wert aktiviert. Pruefen: isls norms fitness
\end{lstlisting}

Wenn SGB \"uber 0.95 steigt: Meilenstein-Highlight.

\begin{lstlisting}[style=bash]
  * Meilenstein: SGB = 0.96
    Nur noch 2 von 50 Dateien benoetigen LLM-Generierung.
    Ziel SGB = 1.0: 2 Norm-Promotions entfernt.
\end{lstlisting}


% ══════════════════════════════════════════════════════════════════════════════
\section{Neue Dateien}
% ══════════════════════════════════════════════════════════════════════════════

\begin{longtable}{lp{10cm}}
\toprule
\textbf{File} & \textbf{Purpose} \\
\midrule
\texttt{isls-gateway/src/timeseries.rs}
  & \texttt{TimeseriesEntry}, \texttt{collect\_snapshot()}, 
    \texttt{append\_entry()}, \texttt{read\_entries()}, 
    Background-Recorder-Task, \texttt{GET /api/timeseries} 
    Handler. \\
\texttt{isls-gateway/src/auto\_evolve.rs}
  & Auto-Evolve Background-Task, Description-Pool, 
    Scrape-after-Forge, Toggle-Handler. \\
\bottomrule
\end{longtable}


% ══════════════════════════════════════════════════════════════════════════════
\section{Modifizierte Dateien}
% ══════════════════════════════════════════════════════════════════════════════

\begin{longtable}{lp{10cm}}
\toprule
\textbf{File} & \textbf{Change} \\
\midrule
\texttt{isls-gateway/src/lib.rs}
  & \texttt{pub mod timeseries; pub mod auto\_evolve;} 
    Timeseries-Recorder und Auto-Evolve-Task spawnen 
    bei Server-Start. Route \texttt{/api/timeseries}. 
    Toggle-Route \texttt{POST /api/auto-evolve/toggle}. 
    Status-Route \texttt{GET /api/auto-evolve/status}. \\
\texttt{isls-gateway/src/static/studio.html}
  & Vierter Modus $\blacklozenge$ im Mode-Strip. 
    Sensorium-Layout mit 6 Karten. SVG-Rendering-Funktionen. 
    Zeitfenster-Umschaltung. Auto-Evolve-Toggle. 
    30-Sekunden auto-refresh. \\
\texttt{isls-cli/src/main.rs}
  & \texttt{--auto-evolve} Flag f\"ur \texttt{serve} 
    Subcommand. \\
\texttt{isls-forge-llm/src/metrics.rs}
  & \texttt{sgb} Feld in \texttt{GenerationMetrics} 
    berechnen und speichern. \\
\bottomrule
\end{longtable}


% ══════════════════════════════════════════════════════════════════════════════
\section{Constraints for Claude Code}
% ══════════════════════════════════════════════════════════════════════════════

\begin{enumerate}
\item \textbf{Regel 13:} Alles muss im Studio sichtbar sein. 
      Kein backend-only Feature. Das Sensorium ist 
      \emph{prim\"ar} ein Studio-Feature.
\item \textbf{Kein Chart-Framework.} Kein D3, kein Chart.js, 
      kein Plotly. Reines SVG, inline generiert. Die Graphen 
      sind einfach genug f\"ur handgeschriebenes SVG.
\item \textbf{Timeseries ist append-only JSONL.} Kein 
      Datenbank-Backend. Kein SQLite. Eine Datei, eine 
      Zeile pro Snapshot, 96 Zeilen pro Tag.
\item \textbf{Auto-Evolve ist default AUS.} Muss explizit 
      aktiviert werden. Kein Verbrauch von Ressourcen 
      ohne Zustimmung.
\item \textbf{Auto-Evolve-Ergebnisse werden nach Scrape 
      gel\"oscht.} Kein Disk-Verbrauch durch generierte 
      Projekte die niemand braucht.
\item \textbf{Sensorium hat KEIN Input-Feld.} Kein Chat. 
      Nur Visualisierungen. Der User schaut, liest, 
      versteht --- tippt nichts.
\item \textbf{Responsive:} Sensorium funktioniert auf 
      Mobile (1 Spalte, Karten stacken vertikal).
\item \textbf{Do not modify} isls-merkaba, isls-reader, 
      isls-code-topo, isls-hypercube, isls-types, isls-chat, 
      isls-norms.
\item Alle bestehenden Tests MUST pass.
\item Branch: \texttt{claude/implement-isls-i5-spec}.
\item Commit format: \texttt{I5/W\{n\}: <description>}.
\end{enumerate}


% ══════════════════════════════════════════════════════════════════════════════
\section{Dependency Graph}
% ══════════════════════════════════════════════════════════════════════════════

\begin{verbatim}
W1 (Timeseries-Persistenz)
  |
  +---> W3 (SVG-Rendering-Engine)
  |       |
  |       v
  +---> W2 (Sensorium-Modus im Studio)
  |
  +---> W4 (Autonomer Evolve-Zyklus)
          |
          v
        W5 (SGB-Tracking + Konvergenz-Alarm)
\end{verbatim}

W1 first (Datengrundlage). W3 parallel (Rendering-Funktionen). 
W2 braucht W1 + W3. W4 unabh\"angig (Backend-Task). 
W5 braucht W1 + W4.

Empfehlung: W1 $\to$ W3 $\to$ W2 $\to$ W4 $\to$ W5.


% ══════════════════════════════════════════════════════════════════════════════
\section{Acceptance Criteria}
% ══════════════════════════════════════════════════════════════════════════════

\begin{longtable}{clp{9cm}}
\toprule
\textbf{\#} & \textbf{Pass} & \textbf{Criterion} \\
\midrule
1 & & \texttt{cargo build --workspace} compiles. \\
2 & & \texttt{cargo test --workspace} passes. \\
3 & & \texttt{timeseries.jsonl} wird alle 15 Minuten geschrieben. \\
4 & & \texttt{GET /api/timeseries?hours=24} liefert Daten. \\
5 & & Studio: $\blacklozenge$ Modus zeigt 6 Karten. \\
6 & & SGB-Trajektorie als Liniengraph sichtbar. \\
7 & & Norm-Wachstum als Liniengraph sichtbar. \\
8 & & Fitness-Verteilung als Balkendiagramm sichtbar. \\
9 & & Scrape-Rate als Liniengraph sichtbar. \\
10 & & Codematrix-Radar als Pentagon sichtbar. \\
11 & & Gen-Topologie als Cluster-Map sichtbar. \\
12 & & Zeitfenster umschaltbar (24h / 7d / 30d). \\
13 & & Auto-Evolve Toggle im Sensorium. \\
14 & & Mobile: Karten stacken vertikal. \\
15 & & Alle bestehenden Tests bestehen. \\
\bottomrule
\end{longtable}


% ══════════════════════════════════════════════════════════════════════════════
\section{Nach I5: Die Roadmap}
% ══════════════════════════════════════════════════════════════════════════════

\begin{longtable}{clp{8cm}}
\toprule
\textbf{Phase} & \textbf{Status} & \textbf{Was} \\
\midrule
D1--D8 & $\checkmark$ & Donnerkeil (Pipeline) \\
M1     & $\checkmark$ & MERKABA Orbitkern (Swarm) \\
S1     & $\checkmark$ & Studio Rebuild \\
I1     & $\checkmark$ & Infogenetik-Kern (5D, Gates, Fitness) \\
I2     & $\checkmark$ & Evolution (Gen-Clustering, Ollama) \\
I3     & $\checkmark$ & Metamorphose (Spectroscopy, Evolve) \\
I4     & $\checkmark$ & Cross-Language Barbara \\
I5     & Spec & Sensorium (Dashboard, Auto-Evolve) \\
\midrule
R3     & Offen & Selbstverbesserung (Subsysteme neu generieren) \\
R4     & Offen & Selbstreproduktion (HDAG-Builder generieren) \\
R5     & Offen & Quine (vollst\"andige Selbstreplikation) \\
\bottomrule
\end{longtable}

Nach I5 sind alle operativen Systeme installiert. R3--R5 sind 
experimentell --- jede Stufe erfordert dass die vorherige 
funktioniert. R3 kann beginnen sobald die SGB \"uber 0.95 
liegt und die Norm-Dichte f\"ur ISLS-eigene Architektur-Patterns 
ges\"attigt ist.

\vfill
\begin{center}
\color{isls-gray}\rule{0.5\textwidth}{0.4pt}\\[0.5em]
{\small End of I5 Specification.\\[0.3em]
Das System sieht sich jetzt selbst.\\
Es misst, es visualisiert, es lernt.\\
Und wenn du schl\"afst, evolvert es.}
\end{center}

\end{document}
